% !TeX root = ../../Book5.tex
% 纲要标题：### B.1 量子群与 Hopf 代数
% 纲要位置：工作记录/计划纲要.md，第 4577 行。
% #### B.1.1 生成关系与余乘法
% #### B.1.2 权模与量子整数
% #### B.1.3 张量次序与非余交换性
\section{量子群与 Hopf 代数}
\label{b5:sec:B01}

本节以 \(U_q(\mathfrak{sl}_2)\) 介绍量子群，验证它的 Hopf 结构，并比较参数对表示的影响。量子半单代数的一般理论及其分类可留待进一步学习。

\subsection{生成关系与余乘法}
\label{b5:subsec:B01:01}

\begin{definition}\label{b5:def:B-quantum-sl2}
设 \(k\) 为特征零域，\(q\in k^\times\)、\(q^2\ne1\)。由生成元 \(E,F,K,K^{-1}\) 及关系
\[
KK^{-1}=K^{-1}K=1,\quad KEK^{-1}=q^2E,\quad KFK^{-1}=q^{-2}F,\quad
EF-FE=\frac{K-K^{-1}}{q-q^{-1}}
\]
定义的含幺结合 \(k\) 代数，记为 \(U_q(\mathfrak{sl}_2)\)，称为相应的量子包络代数。
\end{definition}
\begin{proposition}\label{b5:prop:B-quantum-hopf}
设 \(k\) 为特征零域，\(q\in k^\times\)、\(q^2\ne1\)，\(U=U_q(\mathfrak{sl}_2)\) 为由 \(E,F,K,K^{-1}\) 生成的含幺结合 \(k\) 代数，关系为
\[
KK^{-1}=K^{-1}K=1,\quad KEK^{-1}=q^2E,\quad KFK^{-1}=q^{-2}F,\quad
EF-FE=\frac{K-K^{-1}}{q-q^{-1}}.
\]
则以下公式延伸为代数同态 \(\Delta:U\to U\otimes_kU\)、\(\varepsilon:U\to k\) 和反代数同态 \(S:U\to U\)，并给出 \(U\) 的 Hopf 结构：
\[
\begin{aligned}
\Delta(K)&=K\otimes K,&
\Delta(E)&=E\otimes K+1\otimes E,&
\Delta(F)&=F\otimes1+K^{-1}\otimes F,\\
\varepsilon(K)&=1,&\varepsilon(E)&=0,&\varepsilon(F)&=0,\\
S(K)&=K^{-1},&S(E)&=-EK^{-1},&S(F)&=-KF.
\end{aligned}
\]
\end{proposition}
\begin{proof}
先检查余乘法在自由代数上定义的同态保持关系。关于 \(K\) 的逆及共轭等式逐项由 \(KE=q^2EK\)、\(KF=q^{-2}FK\) 得到。交换子中两项交叉乘积相消，因为
\(EK^{-1}=q^2K^{-1}E\)、\(KF=q^{-2}FK\)。故
\[
\Bigl[\Delta(E),\Delta(F)\Bigr]
=[E,F]\otimes K+K^{-1}\otimes[E,F]
=\frac{K\otimes K-K^{-1}\otimes K^{-1}}{q-q^{-1}},
\]
正是右侧关系的余乘法。余单位也保持关系。对极按反同态延拓；例如
\(S(E)S(K)-q^2S(K)S(E)=0\)
对应 \(KE=q^2EK\) 的反序像；\(F\) 的共轭关系同理由
\(KF=q^{-2}FK\) 给出。剩下的交换子关系为
\[
 S(F)S(E)-S(E)S(F)
 =KFEK^{-1}-EF
 =FE-EF
 =\frac{K^{-1}-K}{q-q^{-1}},
\]
正是 \(S\Bigl((K-K^{-1})/(q-q^{-1})\Bigr)\)；这里 \(K\) 与 \(FE\) 交换。逆元关系也被保持，故三映射在商代数上良定义。
余结合与余单位只需在 \(E,F,K^{\pm1}\) 上核验；例如两种三重余乘法在 \(E\) 上都为
\(E\otimes K\otimes K+1\otimes E\otimes K+1\otimes1\otimes E\)。
对极的左、右等式在 \(E\) 上分别为
\(-EK^{-1}K+E=0\)、\(EK^{-1}-EK^{-1}=0\)，在 \(F\) 上分别为
\(-KF+KF=0\)、\(F-K^{-1}KF=0\)，在 \(K^{\pm1}\) 上都是逆元等式。
还须说明生成元上的核验能推广到乘积。用
\(\Delta(a)=\sum a_{(1)}\otimes a_{(2)}\) 的记号，若左对极等式已对 \(a,b\) 成立，则
\[
\begin{aligned}
\sum S\Bigl(a_{(1)}b_{(1)}\Bigr)a_{(2)}b_{(2)}
&=\sum S\Bigl(b_{(1)}\Bigr)S\Bigl(a_{(1)}\Bigr)a_{(2)}b_{(2)}\\
&=\varepsilon(a)\sum S\Bigl(b_{(1)}\Bigr)b_{(2)}
=\varepsilon(a)\varepsilon(b)1.
\end{aligned}
\]
右等式按相反次序先消去 \(b\) 的两项，得到同一结论。由乘积归纳及线性性，全部对极等式成立。
\end{proof}
这些公式及其版本可对照\cite[Definition 5.6.1; Theorem 5.6.2]{EtingofEtAl2015TensorCategories}。同一量子代数有不同但等价的余乘法惯例，计算张量表示时必须选定一套，不能混合来自两套惯例的公式。

\subsection{权模与量子整数}
\label{b5:subsec:B01:02}

对整数 \(r\)，定义量子整数
\[
[r]_q=\frac{q^r-q^{-r}}{q-q^{-1}}.
\]
对 \(m\ge0\)，在基 \(v_0,\ldots,v_m\) 上规定
\[
Kv_i=q^{m-2i}v_i,\qquad Fv_i=v_{i+1},\qquad
Ev_i=[i]_q[m-i+1]_qv_{i-1},
\]
其中两端之外的向量为零。关于 \(K\) 的两个关系直接由权变化得到；交换子关系等价于
\[
[i+1]_q[m-i]_q-[i]_q[m-i+1]_q=[m-2i]_q,
\]
把各量按定义展开就能核验。所以这些矩阵确实构成一个 \(U_q(\mathfrak{sl}_2)\) 模。
\ProseParagraphBreak
若 \(q\) 不是单位根，则这些 \(K\) 特征值互异，所有内部升算子系数非零。对不变子空间先用 \(K\) 作多项式投影得到某个基向量，再升到 \(v_0\)、降到所有向量，证明该模不可约。上述结论针对这组最高权模型；一般量子群模的分类还须规定相应的表示类别。

\subsection{张量次序与非余交换性}
\label{b5:subsec:B01:03}

在 \(m=1\) 的二维模型中，\(Ev_1=v_0\)、\(Ev_0=0\)、\(Kv_0=qv_0\)。由固定的余乘法，
\[
E(v_1\otimes v_0)=q\,v_0\otimes v_0,\qquad
E(v_0\otimes v_1)=v_0\otimes v_0.
\]
当 \(q\ne1\) 时，普通交换两因子的线性映射不与 \(E\) 交换。因此即使底层空间仍是通常张量积，表示范畴中的交换同构也不再由直接翻转给出。
\ProseParagraphBreak
在合适的有限维表示类别中，可以引入 \(R\) 矩阵来修正交换，并由此研究辫群作用。\(R\) 矩阵是额外结构，须满足相应的相容恒等式；拟三角 Hopf 代数的一般理论见\cite[Section 8.3]{EtingofEtAl2015TensorCategories}。
